% =====================================================================
% VII. Conclusiones
% =====================================================================
\section{Conclusiones}
\label{sec:conclusiones}

Una U-Net implementada desde cero con VV y VH supera al umbral de Otsu en Sen1Floods11 (IoU de
test 0.669 frente a 0.554), sobre todo porque reduce las falsas alarmas. La pérdida combinada BCE
+ Dice y las conexiones de salto son las decisiones que más pesan. Las etiquetas débiles de Otsu
no aportan, porque no contienen información que la red no extraiga del propio SAR.

El modelo generaliza a un evento no visto con agua calma (Bolivia, brecha de 0.8\,\%), pero
colapsa en Lima 2017. Allí el agua urbana y de huaico se aclara en el radar, las laderas en
incidencia rasante imitan al agua y la etiqueta óptica describe la extensión máxima. El análisis
de fallos confirma cuatro mecanismos físicos con pruebas cuantitativas y descarta un quinto.

Añadir Sentinel-2 eleva el IoU de test a 0.848 ($+0.185$) y corrige justo esos fallos, pero la
ganancia proviene casi por completo del óptico y está favorecida por cómo se etiquetaron los
datos. Para un servicio operativo en el Perú, donde las inundaciones llegan con nubes, el
siguiente paso es dar al modelo una dimensión temporal: una imagen SAR previa al evento para
detectar cambios, y la coherencia interferométrica para el agua urbana. Lima es el caso de
prueba natural para ambos.
